\section{Método experimental}
\label{sec:metodo}

\subsection{Plataforma}

\PorFazer{Preencher a partir de \texttt{manifest.json} de cada campanha: modelo e
memória do Raspberry~Pi~5, sistema operativo e \emph{kernel}, versões do OpenSSL,
do Mosquitto e do compilador, \emph{commit} do código; o mesmo para o computador de
secretária e para as máquinas dos sensores e do atacante; rede (cabo ou Wi-Fi);
dissipador; regulador de frequência; se o NTP esteve ativo.}

Tal como planeado, a avaliação tem dois níveis. No \textbf{nível~1}, os
\emph{microbenchmarks} medem o tempo de processador da cadeia de camadas por
mensagem, sem rede, para 16~B, 256~B, 4~KiB e 64~KiB. No \textbf{nível~2}, o
\emph{broker} e o \emph{gateway} correm no Raspberry~Pi e os sensores e o atacante
noutras máquinas da rede local.

\subsection{Configurações}

Foram medidas as doze configurações da \PUm{Tabela das configurações} (K0 a K11),
sem alterações. O nível~1 tem um perfil para K0, K2 a K7, K9 e K11; K1, K8 e K10
só diferem das vizinhas no transporte e nos limites, que só existem no nível~2.

\subsection{Métricas}

A Tabela~\ref{tab:medidas} indica como cada métrica do plano foi medida.

\begin{table}[ht]
  \centering
  \caption{Como cada métrica foi medida.}
  \label{tab:medidas}
  \footnotesize
  \begin{tabularx}{\textwidth}{@{}L{3.8cm}Y@{}}
    \toprule
    \textbf{Métrica} & \textbf{Medição} \\
    \midrule
    Processador por mensagem & Tempo de processador do processo do \emph{gateway}
      na janela, dividido pelas mensagens recebidas (nível~2); tempo de processador
      por mensagem da cadeia completa (nível~1) \\
    Latência p50 e p99 & Do instante de envio, escrito pelo sensor na mensagem, à
      entrega depois da cadeia de camadas. O relógio de cada máquina de sensores é
      corrigido pelo desvio estimado em relação ao do \emph{gateway} \\
    Débito máximo sustentado & Maior taxa entregue com perda inferior a 1\,\% e
      p99 dentro de $L_{\max}$, numa subida de carga por configuração \\
    Octetos por mensagem & Octetos recebidos na ligação TCP do \emph{gateway} por
      mensagem, separados em conteúdo, envelope, MQTT e restante (TLS e controlo) \\
    Memória & Pico de memória residente do \emph{gateway} e do \emph{broker} na
      janela \\
    Estabelecimento de ligação & Tempo e tempo de processador de religações do
      \emph{gateway} (TLS, CONNECT e subscrições), fora da janela de tráfego \\
    Ataques impedidos, $\pi_a(c)$ & Cada ataque contra cada configuração, com
      tráfego legítimo ativo; maioria sobre as repetições \\
    Fuga de informação & Fração das mensagens legíveis por um subscritor de
      \texttt{\#} e por um atacante na ligação do \emph{gateway} \\
    \bottomrule
  \end{tabularx}
\end{table}

A energia, opcional no plano, não foi medida. \PorFazer{Confirmar, ou descrever o
medidor se vier a ser usado.}

\subsection{Ataques}

O atacante do laboratório executa os ataques da \PUm{Tabela dos ataques}. Dois
tipos de ensaio usam-nos de forma diferente:

\begin{itemize}[leftmargin=*]
  \item A \textbf{matriz de ataques} lança cada ataque contra cada configuração
    completa e dá $\pi_a(c)$, de que resulta $S(c)$. Um ataque travado no
    \emph{broker} (por exemplo, um estranho sem credenciais perante a autenticação
    mútua) conta como impedido por essa configuração.
  \item Os \textbf{cenários} isolam uma defesa de cada vez sobre uma base aberta e
    emparelham cada ataque com um \emph{controlo negativo}, a mesma configuração sem
    essa defesa, em que o ataque tem de ter sucesso. Só assim se atribui uma
    rejeição à camada certa.
\end{itemize}

Os ataques internos partem de um dispositivo comprometido: o atacante recebe as
chaves, a palavra-passe e o certificado de um sensor, e o sensor verdadeiro fica
desligado enquanto o atacante o controla. Os ataques na rede usam um intermediário
na ligação entre o \emph{broker} e o \emph{gateway}, que conta o que consegue ler
e, em MQTT sem TLS, altera, redireciona, atrasa ou suprime mensagens, fazendo de
\emph{broker} malicioso.

\subsection{Procedimento}

Cada configuração foi repetida \PorFazer{10} vezes, em ordem aleatória com semente
registada, com 3~s de estabilização e 30~s de medição por execução. Os intervalos
de confiança são de 95\,\%, por \emph{bootstrap} de percentis sobre as execuções.
O custo marginal de uma camada é a diferença de medianas entre as duas
configurações que só diferem nela, e a diferença das diferenças mede a interação
entre a validade temporal e a anti-repetição e entre o TLS e o envelope. A matriz
de ataques correu separada das medições de custo, porque o intermediário na
ligação do \emph{gateway} acrescenta latência.

Cada campanha guarda os dados brutos de todas as execuções, a ordem dos ensaios, a
semente, a configuração completa e um registo do ambiente (código, plataforma,
versões, temperatura e frequência do processador por execução). Todas as tabelas e
figuras deste relatório são geradas a partir desses dados por um único comando
(Anexo~\ref{sec:reproducao}).

\subsection{Parâmetros fixados antes das medições}
\label{sec:parametros}

Os parâmetros do critério de escolha foram fixados antes da campanha: o ficheiro
que os contém é copiado para a pasta da campanha antes da primeira medição e o seu
resumo criptográfico fica no registo do ambiente.

\begin{itemize}[leftmargin=*]
  \item \textbf{Pesos dos ataques.} Cada linha da tabela de ataques da Parte~1
    pesa~1. Duas linhas são medidas por dois ataques cada (a escuta, na rede e no
    \emph{broker}; a injeção e a adulteração), que pesam 0,5 cada.
  \item \textbf{Pesos do custo.} Processador por mensagem 0,4; latência p99 0,4;
    débito máximo, octetos por mensagem, memória do \emph{gateway} e
    estabelecimento de ligação 0,05 cada.
  \item \textbf{Exequibilidade.} $R_{\text{alvo}} = $ \PorFazer{valor}~mensagens/s
    e $L_{\max} = $ \PorFazer{valor}~ms, escolhidos depois de um piloto no
    Raspberry~Pi (\PorFazer{justificar com o cenário de aplicação}) e antes da
    campanha.
\end{itemize}

\subsection{Desvios ao plano}
\label{sec:desvios}

\begin{itemize}[leftmargin=*]
  \item \textbf{Aceleração AES desligada.} O plano previa
    \texttt{OPENSSL\_armcap=0}. Essa variável desliga todas as extensões ARM que o
    OpenSSL usa, incluindo o NEON, de que o ChaCha20-Poly1305 depende. Um
    dispositivo sem as extensões criptográficas tem, em geral, NEON; por isso
    mediu-se também com \texttt{OPENSSL\_armcap=1} (só NEON), e é essa a
    comparação usada para a QP4. Os dois resultados são apresentados.
  \item \textbf{Estabelecimento de ligação.} O tempo de ligação medido varia em
    degraus de cerca de 100~ms, por esperas internas da biblioteca cliente, que
    escondem os poucos milissegundos do TLS. No custo agregado entra o tempo de
    processador por ligação; o tempo decorrido é reportado à parte.
  \item \textbf{Escuta no \emph{broker}.} Na matriz de ataques, este ataque é um
    subscritor de \texttt{\#} sem credenciais. Um \emph{broker} curioso não é
    medido por $S(c)$ (Secção~\ref{sec:ameacas}); é o requisito de
    confidencialidade do caso de uso que o cobre.
  \item \textbf{Dispositivo comprometido.} O sensor verdadeiro fica desligado
    durante o ataque. Com os dois ligados, disputam a sessão MQTT e o ataque fica
    reduzido a uma fração do tempo (Secção~\ref{sec:observacoes}).
  \item \PorFazer{Outros desvios que surjam na campanha (repetições a mais,
    execuções repetidas, energia).}
\end{itemize}
